\chapter{Lista 2.}
\begin{problem}[1]

	Niech G będzie grupą. Udowodnić, że dla dowolnego $g\in G$, $m,n \in \Z$ zachodzą wzory $g^{n+m} = g^{m}g^{n}$ i $\left( g^{m} \right)^{n} = g^{nm}$
\end{problem}
\begin{solution}
	TODO
	Niech G będzie grupą. Ustalmy dowolne $g\in G$.
\begin{itemize}
	\item $g^{n+m} = g^{m}g^{n}$ 

		Rozpatrzymy najpierw przypadek dla liczb naturalnych.

		Ustalmy dowolne $m \in N$.
		Dla $n=0$ mamy
		 \[
			 L = g^{n+m} = g^{0+m} = g^{0}g^{m} = eg^{m} = g^{m}e= g^{m}g^{0} = g^{m}g^{n} = P
		.\] 
		Zatem baza indukcji zachodzi.
	\item  $\left( g^{m} \right)^{n} = g^{nm}$
		Ustalmy teraz dowolne $n>0$ i załóżmy, że $g^{n+m} = g^{m}r^{n}$.
                Wtedy
		\begin{align*}
			L =& g^{(n+1) + m} = g^{(n+m) + 1} = g^{n+m}g^{1}  \overset{\text{z.ind.}}{=} \\ =& g^{n}g^{m}g^{1} = g^{n}g^{m+1} =g^{n}g^{1+m} = g^{n}g^{1}g^{m} = g^{n+1}g^{m} = P 
		.\end{align*}
		Krok indukcyjny zachodzi. Zatem  teza jest prawdziwa dla wszystkich liczb naturalnych.

		Ustalmy teraz $n,m\in \Z$. 
		Mamy cztery możliwości:
		\begin{enumerate}
			\item $n \ge 0$ i $m\ge 0$: to już pokazaliśmy,
			\item $n \ge 0$ i $m<0$:

				Niech $-m' = m$, czyli  $m' \in N^{+}$. Pokażemy przez indukcję względem $m'$, że  $g^{n-m'} = g^{n}g^{-m'}$.

				Dla $m' = 1$ mamy
				 \begin{align*}
					 L & = g^{n-m'} = g^{n-1} = g^{0-1} = g^{-1} = eg^{-1} = g^{0}g^{-1} = P \text{ (dla n = 0)} \\
					 L & = g^{n-m'} = g^{n-1} = \text{z definicji jest to n-1 krotne mnozenie n} = g^{n}g^{-1} = P 
				 .\end{align*}
				 Załóżmy, że dla dowolnego $m' > 0$ zachodzi, wtedy
				 \[
				 L = g^{n-(m'+1)} = g^{(n-m') - 1} = g^{n-m'}g^{-1}= g^{n}g^{-m'}g^{-1} = g^{n}g^{-(m+1)} = P
				 .\] 

				Czyli krok zachodzi.
		        \item $n<0$ i  $m\ge 0$ analogicznie to przypadku powyższego.,
			\item $n<0$ i  $m<0$ 
				Mamy $n = -p$ oraz $m = -q$
				czyli \[
				g ^{n+m} = g^{-p-q} = g^{-1(p+q)} = g^{-1}g^{p+q}
				.\] 
		\end{enumerate}

\end{itemize}
\end{solution}

\begin{problem}[2]

	Niech $\left(A, +  \right)$ będzie grupą przemienną i $m\in \Z$. Udowodnić, że funkcja
	\[
	f: A \rightarrow A,\quad f(a) := ma
	.\] 
	jest homomorfizmem.
\end{problem}
\begin{solution}

	Weźmy dowolne elementy $a_1,a_2 \in A$. Wtedy \[
		f(a_1+a_2) = m(a_1+a_2) = ma_1 + ma_2 = f(a_1) + f(a_2)
	.\] 	
	Zatem f jest homomorfizmem.
\end{solution}
\begin{problem}[3]

	Znaleźć izomorfizm pomiędzy $D_3 i S_3$.
\end{problem}
\begin{solution}

	TODO
\end{solution}
\begin{problem}[4]

	Udowodnić, że $S^{1} \cong SO_2(\R)$.
\end{problem}
\begin{solution}
	
	TODO
\end{solution}
\begin{problem}[5]

	Udowodnić, że:
	\begin{enumerate}[label=\alph*)]
		\item złożenie homomorfizmów jest homomorfizmem,
		\item złożenie izomorfizmów jest izomorfizmem,
		\item funkcja odwrotna do izomorfizmu jest izomorfizmem,
		\item jeśli $G\cong H$ oraz $H \cong N$ to $G \cong N$.
	\end{enumerate}
\end{problem}
\begin{solution}

	Niech  $f: G \rightarrow H$ oraz $g: Q \rightarrow P$ będą dowolnymi homomorfizmami.
	\begin{enumerate}[label=\alph*)]
		\item TODO
		\item TODO
		\item TODO
		\item Ustalmy dowolne grupy $G,H,N$ i załóżmy, że  $G\cong H$ oraz $H \cong N$.
				\begin{center}
					\begin{tikzcd}[row sep= 30pt, column sep=30pt]
						G \arrow[dr, "\Pi"] \arrow[rr, "\Phi", dashrightarrow] & & N \\					     
						& \arrow[ur, "\Psi", rightarrow]H &  
					\end{tikzcd}
				\end{center}
				Złożenie izomorfizmów jest izomorfizmem, zatem między $G$ i $N$ istnieje izomorfizm $\Pi \circ \Psi : G \rightarrow N$
	\end{enumerate}
\end{solution}
\begin{problem}[6]
	
	Pokazać, że grupa $\left( \Q, + \right)$ nie jest skończenie generowana.
\end{problem}
\begin{solution}

	Dana grupa $G$ jest skończenie generowana jeśli spełnia poniższy warunek:
	\[
		\exists_{g_1,g_2,\ldots,g_{n} \in G} \langle g_1,g_2,\ldots,g_{n} \rangle = G   
	.\] 

	\begin{remark}
	
		W grupie Abelowej jaką jest $\left( \Q, + \right) $, $k$-krotne dodawanie zapisujemy jako
		 \[
			 \underbrace{q + q + q + \ldots + q}_{\text{k razy}} = k \cdot n
		.\] 
	\end{remark}	

	Załóżmy, że $\left( \Q, + \right) $ jest skończenie generowana, więc istnieje skończony n-elementowy układ generatorów- są to jakieś liczby wymierne. Jest on postaci \[
	\langle \frac{w_1}{q_1}, \frac{w_2}{q_2},\ldots, \frac{w_{n}}{q_{n}} \rangle 
	.\] 
	Dowolna liczba wymierna, będzie postaci \[
	\frac{q}{\sum_{i=1}^{n} q_{i}}
	.\] 
	gdzie $q$ jest jakąś liczbą całkowita. Niech $p$ będzie największą liczbą pierwszą która dzieli  $\sum_{i=1}^{n} q_i$, wtedy dla liczby wymiernej, która ma w mianowniku liczbę p'- kolejną liczbę pierwszą po p, nie jest możliwe znalezienie takiej kombinacji generatorów, która w mianowniku będzie miała p'.
\end{solution}
\begin{problem}[7]

	Niech G będzie grupą i $g\in G$. Przypomnijmy, że rząd(g)$ := |\langle g \rangle| $, przy czym $\aleph$ jest utożsamione z  $\infty$. Udowodnić, że rząd(g) = min$\{n \in \N_{>0} : g^{n} = e\}$ (przyjmujemy, że min$\left( \O \right) = \infty$). 
\end{problem}
\begin{solution}

	TODO		
\end{solution}
\begin{problem}[8]	

	Niech G będzie grupą, g jej elementem, a n liczbą całkowitą różną od 0.
	\begin{enumerate}[label=\alph*)]
		\item Udowodnić, że rząd $\left( g \right) = \text{rząd}\left( g^{-1} \right)$.
		\item Udowodnić, że $g^{n} = e$ wtedy i tylko wtedy, gdy rząd elementu g jest dzielnikiem n.
	\end{enumerate}
\end{problem}
\begin{solution}

	
\end{solution}
\begin{problem}[9]

	Udowodnić, że podgrupa grupy cyklicznej jest cykliczna.
\end{problem}
\begin{problem}[10]

	Niech H i G będą grupami.


	\begin{enumerate}[label=\alph*)]
		\item Udowodnić, że jeśli $f:G\rightarrow H$ jest homomorfizmem grup i $g\in G$ jest elementem skończonego rzędu, to rząd elementu $f(g)$ też jest skończony i dzieli rząd $g$. Pokazać, że jeśli $f$ jest monomorfizmem to rzędy te są równe. 
		\item Niech G będzie grupą cykliczną generowaną przez element g, a h dowolnym elementem H, który w przypadku gdy rząd(g)$<\infty$ spełnia warunek rząd(h)|rząd(g). Udowodnić że istnieje dokładnie jeden homomorfizm $f:G\rightarrow H$, taki że $f(g) = h$.
	\end{enumerate}
\end{problem}
\begin{solution}
	Ustalmy dowolne grupy $G$ i $H$, oraz element  $g\in G$.
	\begin{enumerate}[label=\alph*)]
		\item Niech $g^{n} = e_G$. Wiemy, że homomorfizm mapuje element neutralny na element neutralny odpowiednich grup. f jest homomorfizmem więc
			\[
				f(g^{n}) = \underbrace{f(g)f(g)\ldots f(g)}_{\text{n razy}}= e_H
			.\] 
			Zatem, rząd $f(g)$ jest skończony i wynosi conajwyżej n. 
			Niech rząd $f(g)$ wynosi $d$. Wtedy $\left( f(g) \right) ^{d} = e_H$. Liczbę $n$ możemy przedstawić jako  $n = k\cdot d + r$, gdzie $r \in [0, d-1]$. Wtedy \[
	f(g)^{n} = f(g)^{k\cdot d +r} =f(g)^{k\cdot d}f(g)^{r} = \left( f(g)^{d} \right)^{k}f(g)^{r} = e_H^{k}f(g)^{r} = f(g)^{r}
			.\] 
			Zatem $f(g)^{r} = e_H$. Ponieważ $d$ jest najmniejszą liczbą naturalną, taką że $f(g)^{d}=e_H$ więc $r = 0$, a co za tym idzie $d|n$.  

			Jeżeli ponadto f jest mnonomorfizmem, to każdy z elementów
			\[
			e_H, f(g), f(g)^{2}, f(g)^{3},\ldots, f(g)^{n-1}
			.\] 
			jest różny, zatem $n = d$. 
		\item Załóżmy, że  $\langle g \rangle = G$. 
	\end{enumerate}

	
\end{solution}
\begin{problem}[11]

	Czy istnieje homomorfizm grup $f:G\rightarrow H$, gdzie:

	\begin{enumerate}[label=\alph*)]
		\item $G= \left( \Z,+ \right), H=\left( \Q,+ \right),f(1) = 7$,
		\item 
	\end{enumerate}
\end{problem}
\begin{solution}

	TODO	
\end{solution}
\begin{problem}[12]

	Znaleźć wszystkie homomorfizmy $f:G\rightarrow H$, gdzie:

	\begin{enumerate}[label=\alph*)]
		\item $G = \left( \Z, + \right), H = \left( \Q, + \right),$
		\item $G = \left( \Z_{15}, +_{15} \right), H = \left( \Z_{4}, +_{4} \right)$
		\item $G = \left( Z_{6} \right), H= \left( \Z_{4}, +_{4} \right),$ 
		\item $G = \Z_{2}\times \Z_{2}, H= \left( \Z_{4}, +_{4} \right) .$
	\end{enumerate}
\end{problem}
\begin{solution}

	TODO	
\end{solution}
\begin{problem}[13]
	Niech G będzie grupą skończoną i niech H będzie niepustym podzbiorem G. $H\le G$ wtedy i tylko wtedy, gdy dla każdych $h_1,h_2\in H$ zachozi $h_1h_2\in H.$
\end{problem}
\begin{solution}

	TODO
\end{solution}
\begin{problem}[14]

	Czy podzbiór H grupy G jest podgrupą G, gdzie:
	\begin{enumerate}[label=\alph*)]
		\item $G = \left( \Z\times\Z\times\Z, +\right), H=\{\left( 2x,5x,3x \right):x \in \Z\}.$
		\item $G = \left(\Z_{8}^{*}, *_{8}\right), H=\{1,3\},$
		\item $G = \left( \Z_{4}, +_{4} \right), H= \{0,3\} $,
		\item $G = \left( \Z_{6}, +_{6} \right), H= \{0,3\},$
		\item G jest grupą izometrii kwadratu, a H składa się ze wszystkich czterech obrotów,
		\item G jest grupą izometrii kwadratu, a H składa się ze wszystkich czterech odbić i identyczności.
	\end{enumerate}
\end{problem}
\begin{solution}

	TODO
\end{solution}
\begin{problem}[15]

	Niech $H_1, H_2$ będą podgrupami grupy G.
	\begin{enumerate}[label=\alph*)]
		\item Pokazać, że $H_1\cup H_2$ nie musi być podgrupą G.
		\item Pokazać, że jeśli $H_1\cup H_2$ jest podgrupą G, to $H_1 \subseteq H_2$ lub $H_2 \subseteq  H_1.$
	\end{enumerate}
\end{problem}
\begin{solution}

	TODO
\end{solution}
